\documentclass[10pt,a4paper]{amsart}
\usepackage[margin=19mm]{geometry}
\usepackage{amsmath,amssymb,mathtools}
\usepackage[scr=rsfs,cal=euler]{mathalfa}
\usepackage{xfrac}
\usepackage[unicode,hidelinks,bookmarks=false]{hyperref}
\newcommand{\ie}{i.e.\ }
\newcommand{\cf}{cf.\ }
\newcommand{\resp}{respectively\ }
\newcommand{\Subsection}{\subsection}
\providecommand{\nobreakdash}{\nobreak\hbox{-}}
\setlength{\emergencystretch}{2em}
\multlinegap0pt
\usepackage{bm}
\usepackage{bbm}
\usepackage{dsfont}
\usepackage{xspace}
\usepackage{cancel}
\usepackage{mathtools}
\usepackage{amsthm}
\makeatletter
\@ifundefined{theorem}{\newtheorem{theorem}{Theorem}}{}
\@ifundefined{lemma}{\newtheorem{lemma}[theorem]{Lemma}}{}
\makeatother
\title[Digraph-defined external difference families and new circula]{Digraph-defined external difference families and new circular external difference families}
\author{Sophie Huczynska, Christopher Jefferson, Struan McCartney}
\date{Source: arXiv:2504.20959v2 (CC BY 4.0)}
\begin{document}
\maketitle

\noindent\emph{Source and licence.} Digraph-defined external difference families and new circular external difference families. Version arXiv:2504.20959v2 (CC BY 4.0), distributed under the Creative Commons Attribution 4.0 International licence, as recorded in the arXiv licence metadata for this exact version and shown on the abstract page https://arxiv.org/abs/2504.20959v2. Full rights notice: licenses/arxiv-2504.20959-CC-BY-4.0.txt.\par\smallskip

\noindent\textit{Editorial scope.} Counted unit: the Theorem whose source label is \texttt{thm:m=4} in the article (source0.tex lines 537--550, source label \texttt{thm:m=4}), delivered together with the article's own complete original proof of it (source0.tex lines 551--606, one \texttt{proof} environment of 56 lines). No mathematics is written, repaired or re-derived: statement and proof are the article's own text line for line. Required context retained uncounted: lem:intdiffs (lines 522--534). Every definition, display and label of those lines is retained unchanged. Omitted from this package and not redistributed: 1--521, 535--536, 607--811. Nothing inside a retained excerpt is omitted or altered: every retained body is delivered exactly as the article's lines are written, so no omission receipt of any kind accompanies this package. Citations: the retained excerpts carry no citation command at all, so no bibliography entry of the article is transcribed and no reference is redistributed. Reference adaptation: none was needed and none was made; every reference inside the delivered unit resolves inside it, so no unresolved reference or citation is produced. Printed slips kept literal: no printed word, symbol, hypothesis or number of the counted statement or of its proof was altered. Layout and class: the article's own class is \texttt{article} with packages including a4, amsmath, amssymb, amsthm, latexsym, color, graphicx, url, tikz, authblk, fullpage, enumitem, xcolor; the wrapper is the lane's pinned \texttt{amsart} editorial wrapper at 10pt on A4, which restates the theorem-like environments the retained text needs and adds the article's own macro definitions that the retained text needs (none), with extra wrapper packages (bm, bbm, dsfont, xspace, cancel, mathtools). The delivered numbering is the unit's own. Assets: no retained span contains a figure environment, a graphics-inclusion command, a PDF-page inclusion, a table or a TikZ picture; no image file is copied into this package, the package declares no asset dependency and no image file is redistributed with it. Licence: version arXiv:2504.20959v2 of this article is distributed under the Creative Commons Attribution 4.0 International licence, as recorded in the arXiv licence metadata for this exact version and shown on the abstract page https://arxiv.org/abs/2504.20959v2. A full rights notice is delivered at licenses/arxiv-2504.20959-CC-BY-4.0.txt. Characters: every retained excerpt is pure ASCII, so the delivered engine, XeLaTeX with its default Latin Modern text font, reports no missing character.
\begin{lemma}\label{lem:intdiffs}
Let $n \in \mathbb{N}$ and $G=\mathbb{Z}_n$.
\begin{itemize}
\item[(i)]  For an interval $I=[0,k]$ in $G$ with $k<n/2$, the maximum multiplicity of an element in $\Delta(I,I)$ is $k+1$ (attained by element $0$).
\item[(ii)] For an interval $I=[0,k]$ in $G$ with $k<n/2$, the maximum multiplicity of an element in $\Delta(I)$ is $k$ (attained by elements $\pm 1$).
\item[(iii)] For a subset $A$ of $G$ and $\alpha \in G$, the following hold:
\begin{itemize}
 \item [1)] $\Delta(A) = \Delta(A+\alpha)$ and $\Delta(A,A) = \Delta(A+\alpha, A+ \alpha)$ 
\item [(2)] $\Delta(\alpha A) = \alpha (\Delta(A))$ and $\Delta(\alpha A, \alpha A) = \alpha (\Delta(A,A))$
\item [(3)] $\Delta(A+\alpha, A) =\Delta(A,A)+\alpha$.
\end{itemize}
\end{itemize}
\end{lemma}

\begin{theorem}\label{thm:m=4}
Let $l \in \mathbb{N}$ and let $d$ be a divisor of $l$.  Define $\mathcal{A}=(A_0,A_1,A_2,A_3)$ to be the following (ordered) collection of sets in $\mathbb{Z}_{4l^2+1}$:
\begin{itemize}
\item $A_0=\{i: 0 \leq i \leq l-1\}$;
\item $A_1=\bigcup_{k=0}^{d-1}\{\frac{l^2(2k+1)}{d}+(i+1)l: 0 \leq i \leq \frac{l}{d}-1\}$;
\item $A_2=A_0+\frac{l^2}{d}=\{\frac{l^2}{d}+i:0 \leq i \leq l-1\}$;
\item $A_3=A_1+2l^2=\bigcup_{k=0}^{d-1}\{\frac{l^2(2(k+d)+1)}{d}+(i+1)l: 0 \leq i \leq \frac{l}{d}-1\}$.
\end{itemize}
Then
\begin{itemize}
\item[(i)] $\mathcal{A}$ is a $(4l^2+1,4,l,1)$-CEDF in $\mathbb{Z}_{4l^2+1}$.
\item[(ii)] For any two distinct proper divisors $d_1, d_2$ of $l$ such that $l \neq d_1 d_2$, the two $(4l^2+1,4,l,1)$-CEDFs obtained in (i) are non-equivalent. In particular, this guarantees at least two non-equivalent $(4l^2+1,4,l,1)$-CEDFs for any composite $l$.
\end{itemize}
\end{theorem}

\begin{proof}
For (i), we order the elements of $\mathbb{Z}_{4l^2+1}$ in the usual way as $0< \cdots <4l^2$.  Since $d(l-1)<l^2$ (as $d|l$), we have that all elements of $A_0$ precede all elements of $A_2$ (which lie between $l^2/d$ and $l^2/d+l-1$), which in turn precede all elements of $A_1$ (which lie between $l^2/d+l$ and $2d(l^2/d)=2l^2$), which in turn precede all elements of $A_3$ (which lie between $(2d+1)l^2/d+l$ and $4l^2$).  All four sets are pairwise disjoint and have no internal repetition of elements.
We now show that the multiset equation $\cup_{i=0}^{3} \Delta(A_{i+1 \mod 4}, A_i)= (\mathbb{Z}_{4l^2+1} \setminus \{0\})$ holds for $\mathcal{A}$.  We will show that $\Delta(A_1,A_0) \cup \Delta(A_0,A_3)=[1,2l^2]$ and $\Delta(A_2,A_1) \cup \Delta(A_3,A_2)=[2l^2+1,4l^2]$.
\begin{align*}
\Delta(A_1,A_0) &= \bigcup_{k=0}^{d-1}\{\frac{l^2(2k+1)}{d}+(i+1)l-j\}: 0 \leq i \leq \frac{l}{d}-1, 0 \leq j \leq l-1 \}\\
&=\bigcup_{k=0}^{d-1}\bigcup_{i=0}^{\frac{l}{d}-1}\bigcup_{j=0}^{l-1}\{\frac{l^2(2k+1)}{d}+(i+1)l -j \}\\
&=\bigcup_{k=0}^{d-1}\bigcup_{i=0}^{\frac{l}{d}-1}[\frac{l^2(2k+1)}{d}+il+1, \frac{l^2(2k+1)}{d}+(i+1)l ]\\
&=\bigcup_{k=0}^{d-1} [\frac{l^2(2k+1)}{d}+1, \frac{l^2(2k+2)}{d}]
\end{align*}
Using the fact that $-2l^2 \equiv 2l^2+1 \mod 4l^2+1$, we have
\begin{align*}
\Delta(A_0,A_3) &=\Delta(A_0,A_1+2l^2)\\
&=-\Delta(A_1,A_0)-2l^2\\
&= 2l^2+1 - \Delta(A_1,A_0)\\
&= 2l^2+1 + \bigcup_{k=0}^{d-1} [-\frac{l^2(2k+2)}{d},-\frac{l^2(2k+1)}{d}-1]\\
&=\bigcup_{k=0}^{d-1} [\frac{l^2(2d-2k-2)}{d}+1, \frac{l^2(2d-2k-1)}{d}]\\
&=\bigcup_{k=0}^{d-1} [\frac{l^2(2k)}{d}+1, \frac{l^2(2k+1)}{d}]
\end{align*}
where in the final step we take $d-1-k$ instead of $k$.

Hence 
$$\Delta(A_1,A_0) \cup \Delta(A_0,A_3) =\bigcup_{k=0}^{d-1} [\frac{l^2(2k)}{d}+1, \frac{l^2(2k+2)}{d}]=[1,2l^2].$$
Next,
\begin{align*}
\Delta(A_2,A_1) &=\Delta(A_0+\frac{l^2}{d},A_1)\\
&=-\Delta(A_1+,A_0)+\frac{l^2}{d}\\
&=-\bigcup_{k=0}^{d-1} [\frac{l^2(2k+1)}{d}+1, \frac{l^2(2k+2)}{d}]+\frac{l^2}{d}\\
&=\bigcup_{k=0}^{d-1} [\frac{l^2(-2k)}{d}-1, \frac{l^2(-2k-1)}{d}]\\
\end{align*}
Now we add $4l^2+1$ to make the differences positive; note this reverses the direction of the union.
\begin{align*}
&=\bigcup_{k=0}^{d-1} [\frac{l^2(4d-2k-1)}{d}+1, \frac{l^2(4d-2k)}{d}]\\
&=\bigcup_{k=0}^{d-1} [\frac{l^2(2(k+d)+1)}{d}+1, \frac{l^2(2(k+d)+2)}{d}]
\end{align*}
where again in the final step we take $d-1-k$ instead of $k$.
\begin{align*}
\Delta(A_3,A_2) &=\Delta(A_1+2l^2,A_0+\frac{l^2}{d})\\
&=\Delta(A_1,A_0)+2l^2-\frac{l^2}{d}\\
&=\bigcup_{k=0}^{d-1} [\frac{l^2(2k+1)}{d}+1, \frac{l^2(2k+2)}{d}]+2l^2-\frac{l^2}{d}\\
&=\bigcup_{k=0}^{d-1} [\frac{l^2 (2(k+d))}{d}+1, \frac{l^2 (2(k+d)+1)}{d}]\\
\end{align*}
Thus
$$\Delta(A_2,A_1) \cup \Delta(A_3,A_2) =\bigcup_{k=0}^{d-1} [\frac{2l^2(k+d)}{d}+1, \frac{2l^2(k+d+1)}{d}]=[2l^2+1,4l^2]$$
which completes the proof of (i).


For part (ii), let $\mathcal{C}^d=(A_0^{d},A_1^{d},A_2^{d},A_3^{d})$ be the CEDF from part (i) corresponding to divisor $d$ of $l$.  Let $d_1,d_2$ be distinct proper divisors of $l$. If $\mathcal{C}^{d_1}$ and $\mathcal{C}^{d_2}$ were equivalent, the sets of $\mathcal{C}^{d_1}$ could be mapped onto those of $\mathcal{C}^{d_2}$ via a mapping which would preserve the list of multiplicities of the internal differences of each set.   Consider $A_0^d=[0,l-1]$ and $A_2^d=[0,l-1]+\frac{l^2}{d}$ in $\mathbb{Z}_{4l^2+1}$; by Lemma \ref{lem:intdiffs}, the maximum multiplicity of an element in $\Delta(A_0^d)$ or $\Delta(A_2^d)$ is $l-1$.  We will show that, for $A_1^d=\cup_{k=0}^{d-1} (\frac{(2k+1)l^2}{d} +l[1,l/d])$ and $A_3^d=2l^2+A_1^d$, the maximum multiplicity of an element in $\Delta(A_1^d)$ or $\Delta(A_3^d)$ is $\mathrm{max}(l-d,l-\frac{l}{d})$. Note that (since we consider only proper divisors $d$ of $l$) this is $l-1$ precisely if $d=1$.  For distinct proper divisors $d_1,d_2$ of $l$ such that $l\neq d_1 d_2$, the maximum multiplicity of an element in $\{\Delta(A_1^{d_1}), \Delta(A_3^{d_1})\}$ is different from that of an element in  $\{\Delta(A_1^{d_2}), \Delta(A_3^{d_2})\}$.   This implies that it is not possible to map $\mathcal{C}^{d_1}$ onto $\mathcal{C}^{d_2}$ as described, and so these CEDFs are not equivalent. For any composite $l$, we can take $d_1=1$ and $d_2$ any prime divisor of $l$ to obtain two non-equivalent examples.

We now prove the above claim. First consider $A_1$.  Define $J=l([0,\frac{l}{d}-1]+1)$ and $B_k=\{\frac{l^2(2k+1)}{d}+J\}$, so that $A_1=\cup_{k=0}^{d-1} B_k$ and $\Delta(A_1,A_1)=\cup_{0 \leq k_1,k_2 \leq d-1} \Delta(B_{k_2},B_{k_1})$. Observe that $\Delta(B_{k_2},B_{k_1})=\frac{2l^2(k_2-k_1)}{d}+\Delta(J,J)$.  By Lemma \ref{lem:intdiffs}, $\Delta(J,J)=l \Delta([0,\frac{l}{d}], [0,\frac{l}{d}])$. So
\begin{equation}\label{eqn:intdiff}
\Delta(A_1,A_1)=\bigcup_{0 \leq k_1,k_2 \leq d-1}\{\frac{2l^2(k_2-k_1)}{d}+\Delta(J,J)\}.
\end{equation}
All elements of $\Delta(J,J)$ lie within $l[-\frac{l}{d}+1,\frac{l}{d}-1]$ and so all elements of $\Delta(B_{k_2},B_{k_1})$ lie within $[(2(k_2-k_1)-1)\frac{l^2}{d}+l, (2(k_2-k_1)+1)\frac{l^2}{d}-l]$.  It is clear that these multisets do not overlap for distinct values of $k_2-k_1$; moreover since the interval corresponding to $k_2-k_1=d-1$ is $[(2d-3)\frac{l^2}{d}+l,(2d-1)\frac{l^2}{d}-l]$ while that corresponding to $k_2-k_1=-(d-1)$ is $[(2d+1)\frac{l^2}{d}+l+1,(2d+3)\frac{l^2}{d}-l+1]$ there is no ``wraparound" modulo $4l^2+1$ in the subtraction table for $\Delta(A_1,A_1)$. In $\Delta(B_{k_2},B_{k_1})$, the maximum multiplicity of a nonzero element is $\frac{l}{d}$ if $k_1 \neq k_2$, and $\frac{l}{d}-1$ if $k_1=k_2$.

In order to determine the maximum multiplicity of an element in $\Delta(A_1)$, we consider the maximum multiplicity of a nonzero element in $\Delta(A_1,A_1)$.  By Equation (\ref{eqn:intdiff}), we must consider the multiset $\{k_2-k_1: 0 \leq k_1,k_2 \leq d-1\}$, i.e. $\Delta([0,d-1],[0,d-1])$. Applying Lemma \ref{lem:intdiffs} to this multiset, element $0$ attains maximum multiplicity $d$ while element $1$ attains multiplicity $d-1$, the maximum multiplicity in $\Delta([0,d-1])$.  From the $d$ multisets $\Delta(B_{k_2,k_1})$ with $k_2-k_1=0$, i.e. $k_1=k_2$, the maximum possible multiplicity for a nonzero element is $d (\frac{l}{d}-1)=l-d$. From the $d-1$ multisets $\Delta(B_{k_2,k_1})$ with $k_2-k_1=1$, the maximum possible multiplicity for a nonzero element is $(d-1)\frac{l}{d}=l-\frac{l}{d}$ (and clearly all other multiplicities of nonzero elements arising from the $k_1 \neq k_2$ case do not exceed this). Hence the maximum occurrence of a nonzero element is $\mathrm{max}(l-d,l-\frac{l}{d})$.
\end{proof}

\end{document}
